\section{Tool-Exposure Taxonomy}
\label{app:tool-exposure-taxonomy}

This appendix records how environment affordances are passed through to the
underlying model for each normalized harness group in
Appendix~\ref{app:dataset-inventory}. The unit here is the normalized harness,
not the raw run label: reasoning-effort variants and model-name suffixes remain
model-configuration metadata, not new tool-exposure mechanisms. The audit uses
the local agent source when available and trace/config metadata where the
source for a historical scaffold is not present in this repository.

\subsection{Mechanism Index}

\begingroup
\footnotesize
\begin{itemize}
\item \textbf{Native tool schema} (2 harnesses): TAU-bench FewShot;
TAU-bench ToolCalling.
\item \textbf{Code-agent tool registry} (4 harnesses): CORE-Agent; HAL
Generalist Agent; HF Open Deep Research; SciCode Tool Calling Agent.
\item \textbf{Text action protocol} (2 harnesses): Claude Code; SWE-Agent.
\item \textbf{Framework-mediated browser control} (3 harnesses):
AssistantBench Browser Agent; Browser-Use; SeeAct.
\item \textbf{Harness-side augmentation} (3 harnesses): ColBench Example
Agent; SAB Self-Debug; USACO Episodic + Semantic.
\item \textbf{No explicit tools} (1 harness): SciCode Zero Shot Agent.
\end{itemize}
\endgroup

\subsection{Exposed-Tool Inventory}

The joinable source for this inventory is
\texttt{irt\_data/harness\_metadata.csv}, in the \texttt{exposed\_tools} and
\texttt{exposed\_tools\_note} columns. The entries below use normalized harness
labels and summarize the model-visible tools or, for harness-side augmentation,
the harness actions that shape the trace without becoming callable model tools.
The derived binary affordance flags, such as code execution, shell access, web
search, browser control, native domain tools, retrieval context, external
evaluation, and self-debug feedback, are stored in
\texttt{irt\_data/harness\_affordances.csv}.

\begingroup
\footnotesize
\begin{itemize}
\item \textbf{AssistantBench Browser Agent}: browser navigation, click, type,
scroll, page extraction, and done/answer actions mediated by
\texttt{browser\_use.Agent}.
\item \textbf{Browser-Use}: browser navigation, click, type, scroll, page
extraction, and done/answer actions from the Browser-Use framework.
\item \textbf{Claude Code}: shell command, file read/search/edit, test-command,
and submit/finish actions through the external CLI text protocol.
\item \textbf{ColBench Example Agent}: harness dialogue and final-answer
submission; no callable environment tools are exposed to the model.
\item \textbf{CORE-Agent}: \texttt{web\_search}, page visit,
\texttt{python\_interpreter}, \texttt{execute\_bash},
\texttt{inspect\_file\_as\_text}, \texttt{edit\_file},
\texttt{file\_content\_search}, \texttt{query\_vision\_language\_model}, and
\texttt{final\_answer}.
\item \textbf{HAL Generalist Agent}: the CORE-style tools
\texttt{web\_search}, page visit, \texttt{python\_interpreter},
\texttt{execute\_bash}, \texttt{inspect\_file\_as\_text}, \texttt{edit\_file},
\texttt{file\_content\_search}, and \texttt{query\_vision\_language\_model},
plus benchmark-dependent adapters such as TAU-bench airline action wrappers,
AppWorld API execution, and ColBench \texttt{ask\_user}/\texttt{finish\_task}.
\item \textbf{HF Open Deep Research}: \texttt{visualizer},
\texttt{inspect\_file\_as\_text}, a managed \texttt{search\_agent},
\texttt{web\_search}, \texttt{visit\_page}, \texttt{page\_up},
\texttt{page\_down}, \texttt{find\_on\_page\_ctrl\_f},
\texttt{find\_next}, and \texttt{find\_archived\_url}.
\item \textbf{SAB Self-Debug}: complete-program generation, harness code
execution, and self-debug feedback. These are harness-side stages, not
model-callable tools.
\item \textbf{SciCode Tool Calling Agent}: \texttt{web\_search},
\texttt{python\_interpreter}, \texttt{wikipedia\_search}, and
\texttt{final\_answer}.
\item \textbf{SciCode Zero Shot Agent}: no explicit tools.
\item \textbf{SeeAct}: browser observation, click, type, select, scroll,
navigate, and stop/answer actions from the SeeAct browser-action protocol.
\item \textbf{SWE-Agent}: \texttt{bash},
\texttt{str\_replace\_editor} view/create/replace/insert/undo operations,
\texttt{filemap}, \texttt{submit}, and diff-state context.
\item \textbf{TAU-bench FewShot}: airline-native tools
\texttt{book\_reservation}, \texttt{calculate},
\texttt{cancel\_reservation}, \texttt{get\_reservation\_details},
\texttt{get\_user\_details}, \texttt{list\_all\_airports},
\texttt{search\_direct\_flight}, \texttt{search\_onestop\_flight},
\texttt{send\_certificate}, \texttt{think},
\texttt{transfer\_to\_human\_agents},
\texttt{update\_reservation\_baggages},
\texttt{update\_reservation\_flights}, and
\texttt{update\_reservation\_passengers}.
\item \textbf{TAU-bench ToolCalling}: the same airline-native tools as
TAU-bench FewShot, passed through the native \texttt{tools\_info} schema.
\item \textbf{USACO Episodic + Semantic}: retrieval prompts, episodic and
semantic examples, and judge execution as harness-side stages; no interactive
callable environment tools are exposed to the model.
\end{itemize}
\endgroup

\subsection{Canonical Tool-Passage Mechanisms}

\begin{description}
\item[Native tool schema.] The benchmark or harness passes OpenAI-style
tool/function schemas to the model API or to a framework that preserves those
schemas as the model-visible contract. The model returns structured tool calls,
and the harness executes them. This category is reserved for cases where the
environment's own schema, such as \texttt{isolated\_env.tools\_info}, is what
the model is asked to call.
\item[Code-agent tool registry.] Python callables are registered as tools in a
code-agent runtime, most often a \texttt{smolagents.CodeAgent}. The model sees a
tool catalogue through the agent runtime and emits executable code or tool-use
steps that call those Python tools. Some registered tools are thin adapters over
benchmark actions, but the model-visible schema is the wrapper function
signature and docstring, not the benchmark's native tool schema. Thus HAL
Generalist on TAU-bench is classified here: its wrappers eventually call
\texttt{isolated\_env.step(Action(...))}, but the original
\texttt{tools\_info} schema is not passed through directly to the model.
\item[Text action protocol.] Tools are described in the prompt as text and the
model is instructed to emit a parseable action, command, or JSON block. The
harness parses that text and maps it to an environment action.
\item[Framework-mediated browser control.] A browser-agent library supplies an
action space such as click, type, scroll, navigation, and page observation. The
benchmark environment is not passed as domain-specific tool schemas; instead,
the model acts through the browser-control framework.
\item[Harness-side augmentation.] Retrieval, self-debugging, code execution, or
evaluation happens outside the model-visible action interface. The model may see
retrieved context or error messages, but it is not given callable environment
tools for that benchmark step.
\item[No explicit tools.] The scaffold is a direct prompt-to-answer or
prompt-to-code generator for the relevant benchmark. Any execution or grading is
outside the model's callable interface.
\end{description}

Harness-side augmentation spans several strengths. At the light end, a harness
may insert retrieved examples or generated context into a prompt before a normal
model call. At the middle, it may run an external environment loop and feed the
model textual observations without giving it callable tools. At the heavy end,
it may execute generated programs, collect runtime errors, and ask the model to
debug in subsequent turns. In all three cases the augmentation can materially
change behavior, but it is not an environment tool interface exposed to the
model.

\subsection{Normalized Harness Labels}

\begingroup
\footnotesize
\begin{itemize}
\item \textbf{AssistantBench Browser Agent}: framework-mediated browser control.
Uses the \texttt{browser\_use.Agent} stack with LangChain chat models and a
browser object; browser affordances are mediated by the browser-use framework
rather than benchmark-specific schemas.
\item \textbf{Browser-Use}: framework-mediated browser control. Historical
online\_mind2web Browser-Use runs are recorded as browser-use scaffolds. Local
source for this exact scaffold is absent, so the label is assigned from trace
configuration and the Browser-Use harness family.
\item \textbf{Claude Code}: text action protocol / external CLI scaffold.
Claude Code runs are software-engineering agents whose environment actions are
exposed through the CLI/scaffold command interface, not through
benchmark-native tool schemas.
\item \textbf{ColBench Example Agent}: harness-side augmentation, medium
strength. The ColBench scaffold calls the model for textual actions in a
human-interaction environment. Environment stepping is mediated by the harness
dialogue loop rather than by callable model tools.
\item \textbf{CORE-Agent}: code-agent tool registry. CORE-Agent is implemented
as a \texttt{smolagents.CodeAgent} with core Python tools such as web/search,
page visit, Python execution, bash execution, file inspection/editing, and
final answer.
\item \textbf{HAL Generalist Agent}: code-agent tool registry with benchmark
adapters. The general path uses \texttt{smolagents.CodeAgent} and
\texttt{CORE\_TOOLS}. Some benchmarks add adapter tools: TAU-bench actions are
hard-coded wrappers around \texttt{isolated\_env.step}, while AppWorld-style
APIs can be converted from API docs into smolagents tools.
\item \textbf{HF Open Deep Research}: code-agent tool registry. The Open Deep
Research scaffold constructs a smolagents research agent with web, browser,
visualizer, and file-inspection tools. Tools are exposed through the smolagents
runtime.
\item \textbf{SAB Self-Debug}: harness-side augmentation, heavy strength. The
ScienceAgentBench self-debug scaffold asks the model to write complete Python
programs, then the harness executes/debugs externally and may feed error
messages back. The execution environment is not exposed as callable model
tools.
\item \textbf{SciCode Tool Calling Agent}: code-agent tool registry. Uses a
smolagents \texttt{CodeAgent} with web search, Python interpreter, Wikipedia
search, and final-answer tools. These are general scaffold tools rather than
SciCode benchmark-native APIs.
\item \textbf{SciCode Zero Shot Agent}: no explicit tools. Generates code from
prompts with direct model calls. There is no callable tool interface exposed to
the model in this scaffold.
\item \textbf{SeeAct}: framework-mediated browser control. Historical
online\_mind2web SeeAct runs are browser-interaction scaffolds. Local source is
not present in this checkout; classification is based on trace/run metadata and
the SeeAct harness family.
\item \textbf{SWE-Agent}: text action protocol / external CLI scaffold.
SWE-Agent exposes repository operations through its command/tool configuration
and command parser. The model interacts with the environment via scaffold
commands rather than benchmark-native tool schemas. The historical raw label
\texttt{My Agent} is normalized into this group because its trace metadata
points to \texttt{agents/SWE-agent-v1.0} filemap/simple-review configurations;
the raw label is retained in \texttt{irt\_data/harness\_metadata.csv} for
possible exclusion or data-cleaning sensitivity checks.
\item \textbf{TAU-bench FewShot}: native tool schema. The TAU-bench few-shot
harness passes \texttt{isolated\_env.tools\_info} and the environment wiki to
\texttt{FewShotToolCallingAgent}; model actions are structured TauBench tool
calls executed by the environment.
\item \textbf{TAU-bench ToolCalling}: native tool schema. The TAU-bench
tool-calling harness passes \texttt{isolated\_env.tools\_info} to
\texttt{ToolCallingAgent}. Optional perturbations modify the schema and
responses but preserve the native tool-call path.
\item \textbf{USACO Episodic + Semantic}: harness-side augmentation, light to
medium strength. The USACO scaffold performs solve and retrieval stages in the
harness. Retrieved examples or context are inserted into prompts; the model is
not given an interactive callable tool interface for the benchmark environment.
\end{itemize}
\endgroup

\subsection{Tool and Affordance Metadata Construction}

The harness metadata is intentionally split into three related layers:
\texttt{tool\_exposure\_mechanism}, \texttt{exposed\_tools}, and binary
\texttt{affordance\_*} flags. The first layer records how tools or
environment actions are passed to the model. The second records the concrete
tool names or harness-side stages associated with each normalized harness. The
third collapses those tool lists into coarse feature flags that can be joined
onto outcome and trace tables.

\paragraph{Harness normalization.}
Raw scaffold labels are first mapped to the normalized harness labels in
Appendix~\ref{app:dataset-inventory}. This step removes provider/model suffixes
and reasoning-setting suffixes from scaffold names, while preserving the raw
label in \texttt{raw\_scaffold}. For example, historical \texttt{My Agent}
runs are normalized to \textbf{SWE-Agent}, with the original label retained for
possible sensitivity analysis. The normalized harness label is then used as the
unit for mechanism labels, exposed-tool inventories, and affordance flags.

\paragraph{Tool extraction.}
For harnesses with local source, exposed tools are extracted from the agent or
benchmark code. Smolagents-based harnesses are read from their
\texttt{CodeAgent} or \texttt{ToolCallingAgent} construction sites and tool
classes. This yields concrete lists such as \texttt{web\_search},
\texttt{python\_interpreter}, \texttt{execute\_bash},
\texttt{inspect\_file\_as\_text}, \texttt{edit\_file}, and
\texttt{file\_content\_search}. TAU-bench tools are taken from the
environment-native \texttt{isolated\_env.tools\_info} path and the observed
airline action names. SWE-Agent tools are read from the filemap/simple-review
configuration bundles, including shell, file-map, editor, submit, and
diff-state actions. When the exact local source is absent, as for some
historical browser scaffolds, the inventory is recorded at the framework-action
level and the caveat is stored in \texttt{exposed\_tools\_note}.

\paragraph{Harness-side stages.}
For harness-side augmentation, \texttt{exposed\_tools} records stages that
shape the trace rather than model-callable tools. Thus SAB Self-Debug records
complete-program generation, harness code execution, and self-debug feedback;
USACO records retrieval/example stages and judge execution; ColBench records
harness dialogue and final-answer submission. These entries are useful for
affordance analysis but should not be interpreted as native or registered model
tools.

\paragraph{Affordance flags.}
The derived \texttt{affordance\_*} columns are stored once per normalized
harness in \texttt{irt\_data/harness\_affordances.csv} and mirrored onto
\texttt{irt\_data/harness\_metadata.csv}. The flags are:
\begin{itemize}
\item \texttt{affordance\_code\_generation}
\item \texttt{affordance\_code\_execution}
\item \texttt{affordance\_shell\_access}
\item \texttt{affordance\_file\_read}
\item \texttt{affordance\_file\_edit}
\item \texttt{affordance\_web\_search}
\item \texttt{affordance\_web\_page\_visit}
\item \texttt{affordance\_browser\_control}
\item \texttt{affordance\_vision}
\item \texttt{affordance\_native\_domain\_tools}
\item \texttt{affordance\_benchmark\_api\_tools}
\item \texttt{affordance\_human\_dialogue}
\item \texttt{affordance\_retrieval\_context}
\item \texttt{affordance\_external\_judge\_or\_eval}
\item \texttt{affordance\_self\_debug\_feedback}
\item \texttt{affordance\_complete\_program\_prompt}
\end{itemize}
A flag is true when that capability is available through the harness or its
evaluation loop, regardless of whether the capability is exposed as a native
tool schema, a code-agent registered function, a text command, or a
harness-side augmentation.

\paragraph{Interpretation.}
These affordance flags intentionally abstract away from the passage mechanism.
For example, Claude Code and SAB Self-Debug both have code execution or
debugging affordances, but Claude Code exposes execution through an interactive
agent scaffold whereas SAB Self-Debug executes generated programs in the
benchmark harness and may feed errors back into later prompts. Conversely,
TAU-bench ToolCalling and HAL Generalist can both execute airline domain
actions, but the former passes native schemas while the latter re-exports some
actions through smolagents wrappers. Analyses should therefore use affordance
flags together with \texttt{tool\_exposure\_mechanism}, not as replacements for
the mechanism taxonomy.

\subsection{Interpretation for Trace Features}

This metadata separates \emph{available affordances} from \emph{model-visible
tool passage}. For example, TAU-bench ToolCalling and HAL Generalist can both
execute \texttt{update\_reservation\_flights}, but the former passes the
TauBench schema directly through \texttt{tools\_info}, whereas the latter
re-exports that action as a hand-written smolagents wrapper. Likewise, a
browser scaffold and a code-agent scaffold may both access the web, but one
does so through page-control actions and the other through search/page tools.
Downstream trace features should therefore include both coarse affordance flags
(for example browser control, code execution, web search, file editing) and this
tool-passage mechanism label.
